\section{Structural Induction}

\begin{frame}
\centering
\Huge\textbf{Structural Induction}
\end{frame}

% --------------------------------------------------
% Question 1
% --------------------------------------------------

\begin{frame}{Question 1}
A programming language defines arithmetic expressions recursively:

\begin{itemize}
\item Every variable is an expression.
\item If $E$ and $F$ are expressions, then $(E+F)$ is an expression.
\end{itemize}

\textbf{Question:} Prove by structural induction that every expression contains an equal number of opening and closing parentheses.
\end{frame}

\begin{frame}{Answer 1}
Let $P(E)$ be the statement:

\textit{Expression $E$ contains the same number of opening and closing parentheses.}

\textbf{Base case:}

A variable such as $x$ contains no parentheses.

Therefore,

$$
0=0.
$$

Hence, $P(E)$ is true for every basic variable expression.
\end{frame}

\begin{frame}{Answer 1}
\textbf{Inductive step:}

Assume $E$ and $F$ are expressions satisfying $P$.

Therefore, each has an equal number of opening and closing parentheses.

The new expression is

$$
(E+F).
$$

This adds one opening parenthesis and one closing parenthesis.

Thus, equality is preserved.
\end{frame}

\begin{frame}{Answer 1}
Therefore, the resulting expression also contains equal numbers of opening and closing parentheses.

Hence, by structural induction,

$$
\boxed{
\text{Every valid expression has equally many opening and closing parentheses.}
}
$$

\end{frame}

% --------------------------------------------------
% Question 2
% --------------------------------------------------

\begin{frame}{Question 2}
Consider a recursively defined set of binary strings:

\begin{itemize}
\item The empty string $\epsilon$ belongs to the set.
\item If $s$ belongs to the set, then $0s1$ also belongs to the set.
\end{itemize}

\textbf{Question:} Prove by structural induction that every string in the set contains the same number of $0$s and $1$s.
\end{frame}

\begin{frame}{Answer 2}
Let $P(s)$ be:

$$
\#0(s)=\#1(s).
$$

\textbf{Base case:}

For the empty string,

$$
s=\epsilon.
$$

It contains no $0$s and no $1$s.

Therefore,

$$
\#0(\epsilon)=\#1(\epsilon)=0.
$$

\end{frame}

\begin{frame}{Answer 2}
\textbf{Inductive step:}

Assume

$$
\#0(s)=\#1(s).
$$

The construction rule creates

$$
0s1.
$$

This adds exactly one $0$ and one $1$.

Therefore,
$$
\#0(0s1)=\#0(s)+1
$$
and
$$
\#1(0s1)=\#1(s)+1.
$$

\end{frame}

\begin{frame}{Answer 2}
Since

$$
\#0(s)=\#1(s),
$$

we obtain

$$
\#0(0s1)=\#1(0s1).
$$

Hence, by structural induction,

$$
\boxed{
\text{Every string in the set contains equally many }0\text{s and }1\text{s}.
}
$$

\end{frame}

% --------------------------------------------------
% Question 3
% --------------------------------------------------

\begin{frame}{Question 3}
A simple expression language is defined recursively:

\begin{itemize}
\item Every integer constant is an expression.
\item If $E$ and $F$ are expressions, then $(E-F)$ is an expression.
\end{itemize}

\textbf{Question:} Prove by structural induction that every expression contains an equal number of opening and closing parentheses.
\end{frame}

\begin{frame}{Answer 3}
Let $P(E)$ denote:

\textit{$E$ contains the same number of opening and closing parentheses.}

\textbf{Base case:}

An integer constant such as $5$ contains no parentheses.

Thus,

$$
0=0.
$$

Therefore, the property holds for every basic expression.
\end{frame}

\begin{frame}{Answer 3}
\textbf{Inductive step:}

Assume $E$ and $F$ satisfy the property.

Construct

$$
(E-F).
$$

The expressions $E$ and $F$ already contain equal numbers of both types of parentheses.

The construction adds exactly one opening and one closing parenthesis.

Therefore, the equality remains true.
\end{frame}

\begin{frame}{Answer 3}
Hence, the new expression satisfies the property.

By structural induction,

$$
\boxed{
\text{every expression contains equally many opening and closing parentheses.}
}
$$

\end{frame}

% --------------------------------------------------
% Question 4
% --------------------------------------------------

\begin{frame}{Question 4}
A programming language defines binary trees recursively:

\begin{itemize}
\item A single node is a binary tree.
\item If $T_1$ and $T_2$ are binary trees, a new tree can be formed with a root and $T_1,T_2$ as its subtrees.
\end{itemize}

\textbf{Question:} Prove by structural induction that every binary tree contains one more node than the number of edges.
\end{frame}

\begin{frame}{Answer 4}
Let

$$
P(T):\quad
\text{number of nodes}
=
\text{number of edges}+1.
$$

\textbf{Base case:}

A single-node tree has

$$
1\text{ node}
$$

and

$$
0\text{ edges}.
$$

Therefore,

$$
1=0+1.
$$

\end{frame}

\begin{frame}{Answer 4}
\textbf{Inductive hypothesis:}

Assume $T_1$ and $T_2$ satisfy

$$
N_1=E_1+1
$$

and

$$
N_2=E_2+1.
$$

Construct a new tree with a root and subtrees $T_1,T_2$.

The number of nodes is

$$
N=1+N_1+N_2.
$$

\end{frame}

\begin{frame}{Answer 4}
The number of edges is

$$
E=2+E_1+E_2,
$$

because the root contributes two edges to the two subtrees.

Using the inductive hypotheses,

$$
N
=
1+(E_1+1)+(E_2+1)
$$
$$
=E_1+E_2+3
$$

and
$$
E+1
=
(E_1+E_2+2)+1
=
E_1+E_2+3.
$$

Therefore,
$$
N=E+1.
$$

\end{frame}

\begin{frame}{Answer 4}
Hence, by structural induction,

$$
\boxed{
\text{Every binary tree has exactly one more node than edges.}
}
$$

\end{frame}

% --------------------------------------------------
% Question 5
% --------------------------------------------------

\begin{frame}{Question 5}
A recursively defined set of expressions is generated as follows:

\begin{itemize}
\item Every variable is an expression.
\item If $E$ is an expression, then $(E)$ is an expression.
\item If $E$ and $F$ are expressions, then $(E+F)$ is an expression.
\end{itemize}

\textbf{Question:} Prove by structural induction that every expression contains at least one variable.
\end{frame}

\begin{frame}{Answer 5}
Let $P(E)$ denote:

\textit{$E$ contains at least one variable.}

\textbf{Base case:}

Every variable is an expression.

For example,

$$
x
$$

contains one variable.

Therefore, $P(E)$ is true for every basic expression.
\end{frame}

\begin{frame}{Answer 5}
\textbf{Inductive step 1:}

Assume $E$ contains at least one variable.

The construction rule produces

$$
(E).
$$

Parentheses do not remove the variable from $E$.

Therefore, $(E)$ still contains at least one variable.
\end{frame}

\begin{frame}{Answer 5}
\textbf{Inductive step 2:}

Assume $E$ and $F$ each contain at least one variable.

The construction

$$
(E+F)
$$

contains all variables from both $E$ and $F$.

Therefore, it contains at least one variable.

Hence, by structural induction,

$$
\boxed{
\text{Every expression contains at least one variable.}
}
$$

\end{frame}

\begin{frame}
\centering
\vfill
{\Huge\textbf{Thank You}}
\vfill
\end{frame}
